\section{What to do now}\label{sec:plan}

\subsection{What can be claimed today}

Five things, and nothing more.

\begin{enumerate}
\item Inside a topic a tool is worth the time it saves.
\item Insiders' premium over the frontier equals the newcomer's cost of catching up (theory).
\item Growth on arXiv in 2018--2025 opened new topics. This holds qualitatively; the numbers cannot yet be reproduced from the repository and their standard errors are wrong.
\item Papers per author rose by 11--14\% in January--September 2026, after seven flat years.
\item Every equilibrium welfare effect in the model is about 1\% or less.
\end{enumerate}

What cannot be claimed today: that the rule changes identify the model's margins; that a production tool is worth 1.3 to 7.2 times a reading tool; any ranking of tools; that topics are anywhere near a speed limit; that adopters of language models produce 7\% more; that the production shock so far is small.

\subsection{Risks that remain if everything goes well}

\begin{enumerate}
\item \textbf{The test shows that insiders stay, not that the cost of entry is reading.} Only the stock gradient speaks to the second.
\item \textbf{The stakes inside the model are small.} A referee can accept every result and ask why it matters. The answer has to be about quantities: where the output of faster researchers lands, and how fast science spreads it. It cannot be about welfare.
\item \textbf{The shock is young.} Nine months, still building, in a year of changed wording and three rule changes. The most likely outcome is half a result: the extra papers stay home, and the response of entry is unclear.
\item \textbf{Priority.} The 2026 fact will not stay unreported. The paper already cites the two closest studies, on output of adopters of language models and on artificial intelligence narrowing the focus of science; Fact~B and ``speed thickens'' must be positioned against both.
\item \textbf{The odds.} The reviewers put the chance at a top-five journal below one in ten in the best case, and at 2--4\% for the most likely outcome, with a strong field journal or a second-tier general journal as the likely home. These are judgements. What would move them is a reading of 2027 that shows the predicted adjustment. More theory would not.
\item \textbf{A fourth headline.} The draft has changed its headline three times ahead of its evidence. The cure is the order of work: evidence first, introduction last.
\end{enumerate}

\subsection{The first month}

\begin{center}\small
\begin{tabularx}{\textwidth}{@{}P{0.85in}LL@{}}
\toprule
When & The empirical side & The theory side \\
\midrule
Days 1--2 & \multicolumn{2}{p{5.0in}@{}}{Freeze the scope on one page: the working claim, the paper to be written under each outcome, and what is out for good. Move the speed limit, the life cycles, the incidence theorem and the rule-change designs to companion folders. No new theorems. No introduction.} \\
Days 2--6 & Write up the 2026 rise as Fact~B from author names. No text measure is needed. & Derive the short-run prediction for a shock that builds up over a year, from one law of motion with the stock and the catch-up delay as states, for entry priced on the flow and on the stock. \\
End of week 1 & \multicolumn{2}{p{5.0in}@{}}{Deposit the dated memo: outcomes, samples, radii, cell sizes, the construction of the founders' share, thresholds for ``supports'' and ``refutes'', and a power calculation on 2019--25. Nothing about density or entry in 2026 is looked at before this.} \\
Weeks 1--2 & Embed abstracts for 2018 to October 2026. Build flow density, the benchmark, the style floor, the cells and the author-by-cell panel. Validate against the word-based measure on 2019--22. Measure turnover. & Write Appendix~B.8 as one proposition with a proof, with the frictionless comparison. Fix one convention for how topics are born. \\
Weeks 2--3 & \textbf{Gate 1.} Repair Fact~1 on its own panel and build the version on arXiv categories. Stop if the elasticity is not below the benchmark. & Repair Lemma~4, and add a test that computes output from member counts. \\
Weeks 3--4 & \textbf{Gate 2.} Tests A, B and C in the registered order, with the placebo year, the lead test and the versions that use no text. Report entry by authors already on arXiv whatever its sign. & As a hedge: the founders' share in high-energy physics, and newcomers' share against older stock. \\
Week 5 & \multicolumn{2}{p{5.0in}@{}}{Decision meeting. Read the results against the memo. Choose the headline, the wording of the mechanism (``reading'' or ``a sunk entry cost''), and whether tools are a section or a second paper. Only then begin the rewrite.} \\
\bottomrule
\end{tabularx}
\end{center}

\subsection{Logistics}

\begin{itemize}
\item \textbf{The panel behind Facts 1--4 is not in the repository.} Either bring it and its scripts in, or make the arXiv-category version the primary estimate and cite the other as a check. Decide before the facts section is written.
\item \textbf{Sentence embeddings cannot be run in the environment where the code now lives} (no deep-learning library, four cores). The first job of week~1 is a machine that can.
\item \textbf{Reference lists for 2025--26 in high-energy physics} have to be fetched for the text-free check; the script exists (\texttt{empirics/inspire\_fetch.py}).
\item \textbf{The arXiv snapshot} ends on 15 October 2026 and should be refreshed.
\end{itemize}

\subsection{Decisions for the authors}

\begin{enumerate}
\item \textbf{Accept the two gates before any rewriting?} This is empirical work, and it comes before the theory repairs that depend on its outcome.
\item \textbf{Where does Fact~1 live?} On the original panel, brought into the repository, or on a version that runs from it.
\item \textbf{What is retired now, whatever the tests show?} The recommendation: the title, the speed limit as a result, the range 1.3--7.2, the ranking of tools, the 7\%, and the sentence about identifying margins.
\item \textbf{Circulate the 2026 fact early?} One reviewer's advice was to post the fact and the registered predictions as a short working paper now, because the fact will not stay unreported, and to submit the full paper after a first reading of 2027.
\item \textbf{One paper or two?} The recommendation is one now (scale versus speed, or the fallback), with incidence as a second paper once a shock on the joiners' side has outcomes.
\end{enumerate}

\begin{keybox}[The plan in four sentences]
Do not rewrite yet.
Repair the main fact and make it reproducible; derive what the model predicts for 2026 and write it down; then look at where the extra papers of 2026 landed.
Let that choose between \emph{scale versus speed} and the fallback on the entry margin and incidence.
Whatever the outcome, the speed limit leaves the title, and the introduction is written last.
\end{keybox}
